\paragraph{The non-associativity boundary}
The intended language can reject a chain such as $a\leftrightarrow b\leftrightarrow c$ unless it is parenthesized.
The pinned FsYacc 11.3.0 implementation shifts on the equal-precedence non-associative case rather than inserting
an error action. Consequently this edition factors the equivalence production to admit at most one unparenthesized
equivalence at each formula level. The grammar, rather than an unsupported declaration assumption, enforces the language.
@@snippet:lalr-logic@@

\href{../../Chapter-10/Propositional-Logic.ipynb}{\nolinkurl{Chapter-10/Propositional-Logic.ipynb}}
checks precedence and the rejection of chained equivalence.
For a lookahead conflict, compare these two actual specifications:
@@snippet:lalr-lookahead@@
@@snippet:lalr-lookahead-solved@@
The transformation delays reduction until the distinguishing token is available.
The associated \href{../../Chapter-10/03-Look-Ahead.ipynb}{\nolinkurl{03-Look-Ahead.ipynb}}
and \href{../../Chapter-10/04-Look-Ahead-Solved.ipynb}{\nolinkurl{04-Look-Ahead-Solved.ipynb}}
retain the conflicting and resolved input examples.
